\chapter{总结与展望}\label{chap:conclusion}

\section{全文总结}

随着自主移动设备对高精度、高实时性空间感知需求的日益增长，视觉 SLAM 技术在资源受限环境下的落地应用成为了研究的热点。本文针对传统视觉 SLAM 后端优化计算复杂度高、嵌入式 CPU 处理性能不足的问题，提出并实现了一种基于 FPGA 的软硬协同并行加速方案。全文的主要研究工作总结如下：

\begin{enumerate}
    \item \textbf{软硬协同架构设计}：本文深入分析了视觉惯性里程计（VIO）的计算特征，基于异构 SoC 平台设计了合理的任务划分机制。通过将雅可比矩阵计算与舒尔补消元等计算密集型任务卸载至 FPGA 逻辑端，极大地释放了主控 ARM 处理器的负载。
    \item \textbf{高效硬件加速算子开发}：设计并实现了支持 5 级深度流水线的雅可比计算模块，针对 $3 \times 3$ 块矩阵开发了专用的并行化求逆器。通过定点化精度优化，在保证硬件资源利用率的同时，维持了系统的定位精度。
    \item \textbf{稀疏性利用与功耗控制}：创新性地引入了 8 位稀疏观测位图（SOB）编码协议，并据此设计了动态门控机制。该方案有效利用了后端优化中的观测稀疏性，通过状态级旁路和时钟门控技术，不仅提升了计算吞吐率，还显著降低了硬件功耗。
    \item \textbf{全系统集成与定量评价}：在 Zynq UltraScale+ 平台上完成了全系统集成，并利用 EuRoC 等数据集进行了实测。实验结果表明，所设计的加速器相比 ARM CPU 实现了约 3.85 倍的加速比，并能够以 25~Hz 以上的稳定频率运行。
\end{enumerate}

\section{主要创新点}

本文的主要创新点归纳如下：
\begin{enumerate}
    \item \textbf{提出了一种自适应的稀疏观测编码（SOB）机制}，实现了算法理论的稀疏性到硬件逻辑执行流的直接映射，为低功耗 SLAM 硬件设计提供了新思路。
    \item \textbf{构建了面向 SchurVINS 框架的紧凑型存储管理方案}，通过隐式索引和 BRAM 乒乓缓存策略，解决了片上存储资源受限与大规模矩阵并行运算之间的矛盾。
    \item \textbf{实现了一套高并行度的舒尔消元核心}，通过专用算子的深度优化，打破了后端优化在大规模特征点场景下的计算瓶颈。
\end{enumerate}

\section{后续工作展望}

尽管本文在基于 FPGA 的 SLAM 后端加速方面取得了一定成果，但仍存在进一步优化的空间：
\begin{enumerate}
    \item \textbf{支持更复杂的相机模型}：目前系统主要针对针孔模型进行了优化。未来可以探索支持鱼眼相机（Fisheye）或全景相机模型的通用硬件雅可比引擎。
    \item \textbf{全硬件前端集成}：本文的前端追踪仍在 ARM 端运行。后续研究可尝试将特征追踪与光流计算完全整合进 PL 端，实现整套 VIO 系统的全硬件单芯片加速，以获得更高的同步精度。
    \item \textbf{引入机器学习加速器}：结合近年来深度学习在视觉里程计中的应用，探索将传统的几何 SLAM 加速器与神经网络加速器（NPU）相结合，构建语义级的软硬协同感知系统。
\end{enumerate}
